\subsectionaddtolist{نسخه اول}
آسیب‌پذیری این نسخه ناشی از عدم وجود جداکننده (\lr{Delimiter}) مناسب بین مقادیر الحاق شده است. فرمول ساخت \lr{MAC} به صورت زیر است:
$$MAC = H(K \parallel \text{\lr{player\_id}} \parallel \text{\lr{gold\_amount}})$$

با توجه به فرض مسئله، حساب کاربری با نام \lr{eve1} درخواستی برای دریافت $0$ سکه فرستاده و \lr{MAC} معتبر آن را دریافت کرده است. بنابراین مقداری که هش می‌شود برابر است با:
$$K \parallel \text{"\lr{eve1}"} \parallel \text{"0"} = K \parallel \text{"\lr{eve10}"}$$

حال اگر هکر بخواهد برای اکانت اصلی خود یعنی \lr{eve} مقدار $10$ سکه دریافت کند، مقدار ورودی هش سرور به صورت زیر خواهد بود:
$$K \parallel \text{"\lr{eve}"} \parallel \text{"10"} = K \parallel \text{"\lr{eve10}"}$$

همان‌طور که مشاهده می‌شود، هر دو رشته پس از الحاق کاملا یکسان تجمیع می‌شوند ($K \parallel \text{\lr{eve10}}$). در نتیجه، \lr{MAC} تولید شده برای درخواست اول، برای درخواست دوم نیز کاملا معتبر است. 

\textbf{نحوه جعل:} هکر کافی است دقیقا همان \lr{MAC} دریافت شده برای (\lr{eve1}، $0$ سکه) را بردارد و همراه با درخواست جعلی خود یعنی (\lr{eve}، $10$ سکه) به سرور ارسال کند. سرور به دلیل یکسان بودن حاصل الحاق، $MAC$ را معتبر تشخیص می‌دهد.


\subsectionaddtolist{نسخه دوم}
در این نسخه فرمول به صورت زیر تغییر یافته است:
$$MAC = H(K) \oplus H(\text{\lr{player\_id}} \parallel \text{":"} \parallel \text{\lr{gold\_amount}})$$

مشکل اصلی این فرمول استفاده از عملگر \lr{XOR} است. این عملگر خاصیت جابه‌جایی و بازگشتی دارد که به هکر اجازه می‌دهد مقدار ثابت $H(K)$ را کاملا حذف یا ایزوله کند.

ما درخواست معتبر برای کاربری \lr{eve} و مقدار $5$ سکه را شنود کرده‌ایم و $MAC_{\text{\lr{old}}}$ را در دست داریم:
$$MAC_{\text{\lr{old}}} = H(K) \oplus H(\text{"\lr{eve}:5"})$$

از آنجا که مقدار $H(\text{"\lr{eve}:5"})$ بدون نیاز به کلید مخفی و صرفا با داشتن تابع هش قابل محاسبه است، هکر می‌تواند با اعمال \lr{XOR} این مقدار را از $MAC_{\text{\lr{old}}}$ حذف کرده و مقدار مستقل $H(K)$ را استخراج کند:
$$H(K) = MAC_{\text{\lr{old}}} \oplus H(\text{"\lr{eve}:5"})$$

حال هکر می‌خواهد درخواست جدیدی برای دریافت $9999$ سکه جعل کند. فرمول \lr{MAC} جدید برابر است با:
$$MAC_{\text{\lr{new}}} = H(K) \oplus H(\text{"\lr{eve}:9999"})$$

با جایگذاری مقدار استخراج‌شده‌ی $H(K)$، \lr{MAC} جعلی جدید به سادگی به دست می‌آید:
$$MAC_{\text{\lr{new}}} = \left( MAC_{\text{\lr{old}}} \oplus H(\text{"\lr{eve}:5"}) \right) \oplus H(\text{"\lr{eve}:9999"})$$


\subsectionaddtolist{نسخه سوم}
کد پیاده‌سازی شده در سرور دچار آسیب‌پذیری {(\lr{Timing Attack})} است. تابع با استفاده از یک حلقه \texttt{for}، کاراکترهای دو \lr{MAC} را تک‌به‌تک از چپ به راست مقایسه می‌کند و به محض دیدن اولین عدم تطابق، مقدار \texttt{False} را برمی‌گرداند.

{روند استخراج \lr{MAC} توسط هکر:}
\begin{enumerate}
    \item {حدس کاراکتر اول:} هکر \lr{MAC}‌هایی ارسال می‌کند که کاراکتر اول آن‌ها متفاوت است. اگر کاراکتر اول درست باشد، حلقه یک بار بیشتر اجرا شده و زمان پاسخ‌دهی سرور به اندازه یک پالس کوچک (اما قابل اندازه‌گیری) طولانی‌تر می‌شود. هکر با تحلیل آماری زمان پاسخ‌ها، کاراکتر اول \lr{MAC} صحیح را پیدا می‌کند.
    \item {حدس کاراکترهای بعدی:} پس از تثبیت کاراکتر اول، هکر همین روند را برای کاراکتر دوم، سوم و... تکرار می‌کند تا زمانی که زمان پاسخ‌دهی به حداکثر خود برسد (یعنی کل \lr{MAC} شده و حلقه تا انتها اجرا شود).
\end{enumerate}


\subsectionaddtolist{راه حل نهایی}

\begin{enumerate}
    \item {ساختار استاندارد جایگزین:} به جای روش‌های ابداعی، توسعه‌دهندگان باید از ساختار استاندارد {\lr{HMAC (Hash-based Message Authentication Code)}} استفاده کنند.
    
    این ساختار برای جلوگیری از حملاتی مثل \lr{Length Extension} و آسیب‌پذیری‌های الحاق خام، کلید مخفی را در دو مرحله مجزا (با استفاده از دو پد داخلی و خارجی \texttt{\lr{ipad}} و \texttt{\lr{opad}}) با پیام ترکیب کرده و هش می‌کند:
    $$HMAC(K, m) = H\Big( (K \oplus \text{\lr{opad}}) \parallel H\big((K \oplus \text{\lr{ipad}}) \parallel m\big) \Big)$$

    \item {رفع مشکل امنیتی نسخه سوم:} برای جلوگیری از حمله زمان‌بندی، مقایسه دو رشته باید در {زمان ثابت} انجام شود؛ یعنی فارغ از اینکه خطا در کدام کاراکتر رخ داده است، کل رشته همیشه تا انتها بررسی شود.
\end{enumerate}